%% ============================================================
%%  SOLUCIÓN — Ejercicio 3.1: Tabla de estadísticas descriptivas
%%  Clase 1 — Después de diapositiva "Tablas profesionales con booktabs"
%%  Compilar: F5 (pdflatex)
%% ============================================================

\documentclass[12pt, a4paper]{article}
\usepackage[utf8]{inputenc}
\usepackage[T1]{fontenc}
\usepackage[spanish]{babel}
\usepackage{booktabs}

\title{Estadísticas descriptivas de la muestra}
\author{Tu Nombre}
\date{\today}

\begin{document}
\maketitle

\section{Análisis descriptivo}

La Tabla \ref{tab:descriptivas} presenta las principales estadísticas de 
nuestra muestra.

\begin{table}[htbp]
  \centering
  \caption{Estadísticas descriptivas de la muestra}
  \label{tab:descriptivas}
  \begin{tabular}{lrrrr}
    \toprule
    Variable       & Media  & D.E.  & Mín.  & Máx.  \\
    \midrule
    Salario (Gs.) & 2.450.000 & 1.120.000 & 550.000 & 9.800.000 \\
    Años de educación & 12.3 & 3.4  & 0     & 20    \\
    Experiencia (años) & 8.7 & 5.1  & 0     & 40    \\
    Mujer (=1)    & 0.48  & 0.50  & 0     & 1     \\
    \midrule
    Observaciones & \multicolumn{4}{c}{1.200} \\
    \bottomrule
  \end{tabular}
  \begin{minipage}{\linewidth}
    \footnotesize \textit{Fuente:} EPH 2023. Elaboración propia.
  \end{minipage}
\end{table}

\subsection*{Respuestas a las preguntas}

\textbf{¿Qué hace \texttt{\textbackslash multicolumn\{4\}\{c\}\{1.200\}}?}

Este comando une 4 columnas en una sola (``span'') y centra el contenido 
con alineamiento \texttt{c} (center). Es útil para etiquetas que aplican 
a múltiples columnas.

\textbf{¿Qué cambió al alinear la columna ``Variable'' a izquierda?}

La columna 1 pasó de alineamiento \texttt{c} (center, por defecto) a 
\texttt{l} (left), lo que hace más legible la tabla porque los nombres 
se alinean a la izquierda.

\end{document}
